\documentclass[12pt]{article}

\usepackage{amsmath}
\usepackage{german}
\usepackage{graphicx}
\DeclareMathOperator\sgn{sgn}

%\graphicspath{{/home/ritz/Teaching/fcn/fig/}}

%----------------- layout ----------------------------------------

\oddsidemargin=0in
\textwidth=6in
\topmargin=-0.7in
\textheight=9.5in
\baselineskip=16pt
\font\mittel=cmr17
\font\roman=cmr12
\font\klein=cmr9
\font\mgross=cmr10
\newfont{\CAP}{cmcsc10 scaled 900}
\pagestyle{empty}
\parindent=0mm

%----------------- macros ----------------------------------------
\newcommand{\conc}[1]{[{\rm C}]_{\text{#1}}}
\newcommand{\units}[1]{\text{#1}}
\newcommand{\K}{\text{K}}
\newcommand{\Na}{\text{Na}}
\newcommand{\Cl}{\text{Cl}}
\newcommand{\Leak}{\text{leak}}
\newcommand{\water}{\text{H}_2\text{O}}
\newcommand{\tderiv}[1]{\frac{\mathrm{d}#1}{\mathrm{d}t}}
%-----------------------------------------------------------------


\begin{document}

\parbox{2cm}{
\includegraphics[width=1.8cm]{bccnlogo.pdf}
}
\parbox{11cm}{
\begin{center}
\large HUMBOLDT-UNIVERSIT"AT \hskip 0.1 cm ZU \hskip 0.1 cm BERLIN
\vskip 0.1 true cm
\mgross BERNSTEIN CENTER FOR\\ COMPUTATIONAL NEUROSCIENCE
\end{center}
}
\parbox{2cm}
{
\hfill
\includegraphics[width=1.8cm]{hublogo.pdf}
}

\vskip 0.6 true cm

\leftline{\CAP Humboldt-Universit"at zu Berlin 
\hfill Phone: 030/2093-9110}
\leftline{\CAP Philippstr. 13 House 6
              \hfill Fax: 030/2093-6771}
\leftline{\CAP \hfill webpage: http://www.bccn-berlin.de/}


\vskip 0.8 true cm
\centerline{\bf Models of Neural Systems, WS 2010/11}
\centerline{\bf Project 6: Complex cells tuned to disparity}
\centerline{Project presentation and report submission: February, 15th, 2011}

\vskip 0.8 true cm

\noindent {\bf Background}

The visual cortex contains cells sensitive to binocular disparity i.e.
difference in image location of an object seen by the left and right
eye, resulting from the eyes' horizontal separation. The brain uses
binocular disparity to extract the depth information for stereotopic
vision. In the project, you will model receptive fields of such cells
and analyse their responses to gratings and natural images.


\vskip 0.8 true cm
\noindent {\bf Problems}


\begin{enumerate}
    \item \label{simplecell} Construct a model simple cell that sums inputs from both
        eyes. The receptive field of each eye is described by a  Gabor
        wavelet:
        $$\rho(x,y)=\cos(f x' -
        \phi)\exp\left(-\frac{x'^2+y'^2}{2\sigma^2}\right),$$
        where $x'=x \cos\theta+y\sin\theta$;
        $y'=y\cos\theta-x\sin\theta$, $f$ is the cell's preferred
        spatial frequency, $\theta$ is its preferred orientation,
        $\phi$ is its phase and $\sigma$ describes the size of the receptive
        field.
    \item Calculate the response of a simple cell to gratings:
        $$I_{L,R}(x,y)=\cos(Kx+\Phi_{L,R}),$$
        where $I_L$ and $I_R$ are the inputs to the left and right
        eyes respectively; $\Phi_L$, $\Phi_R$ are phase shifts of the
        gratings and $K$ is the spatial frequency. Take
        $\Phi_L=\Phi_R=\Phi$, $f=K=2$, $\sigma=1$.  Plot the tuning curve
        as a function of $\Phi$.
    \item Construct a complex cell that sums squares of responses of
        two simple cells: $L_C =
        L_1^2+L_2^2$. The receptive fields of both simple cells should
        be similar but differ in the phase (for example $\phi_1=0$,
        $\phi_2=\pi/2$).
    \item Show that the complex cell is tuned to disparity in the
        images, $\Phi_L-\Phi_R$, but is independent of the absolute
        spatial phase of the grating, $\Phi_L+\Phi_R$. Plot the
        response tuning curve as a function of disparity. 
    \item Use the above model, estimate the disparity maps in natural
        stereoscopic images. A stereoscopic image is a pair of two images of the same scene, taken from slightly different angles. To this end, calculate the responses of
        complex cells with receptive fields  centered  at each pixel of the image 
        (\textit{Hint:} you can use spatial convolution). Plot obtained disparity
        maps for different preferred frequencies $f$ and preferred
        orientations $\theta$. You can download some stereoscopic
        images from

	\texttt{http://vision.middlebury.edu/stereo/}
        
\end{enumerate}

\begin{thebibliography}{7}
    \bibitem{dayan2001} P Dayan, LF Abbott, ``Theoretical neuroscience'',
        MIT Press 2001
    \bibitem{deAngelis1991} GC DeAngelis, I Ohzawa, RD Freeman,
        ``Depth is encoded in the visual cortex by a specialized
        receptive field structure'', Nature 352: 156-159 (1991)
\end{thebibliography}

\vfill
\centerline{\CAP Contact}
\CAP

\begin{tabular}{lll}
\underline{Matthias Guggenmos} & Phone: 450-517131 & Email: matthias.guggenmos(at)bccn-berlin.de \\
Urs Bergmann & Phone: 2093-8924 & Email: urs.bergmann(at)biologie.hu-berlin.de \\
Bartosz Telenczuk & Phone: 2093-8838 & Email: b.telenczuk(at)biologie.hu-berlin.de \\
Richard Kempter & Phone: 2093-8925 & Email: r.kempter(at)biologie.hu-berlin.de \\
\end{tabular}
\end{document}


\end{document}
 
